Gの写真F/Dは各42移動で非supervisor保存ファイルが一致したが，公式成否は未達／成功，通知は０／１件だった．まずGの写真F/Bを別環境で再現し，元Dと同じ空引数の完了通知なし／１回追加の４条件を固定した．
次にNの全12条件のうち，成否によらず完了API未送信の全７状態（既知３課題）を選び，同じ２条件の計14ケースを固定した．Nの通知は元写真Bで観測した\texttt{status="success"}だけとした．全ケースで送信済み公開APIの返却と介入前の全13保存ファイルを元状態と照合し，無介入は元の成否を再現した．
\begin{center}\small 表4\quad 完了通知だけの介入（全18ケース）\par\smallskip
\begin{tabular}{lll}\hline 再現した元状態 & 通知なし & 通知１回\\\hline
G 写真 F & 未達 & 成功\\
G 写真 B & 未達 & 未達\\
N 写真 D & 未達 & 成功\\
N 精算 F & 未達 & 未達\\
N 精算 D & 未達 & 未達\\
N 精算 B & 未達 & 未達\\
N 書出し F & 未達 & 未達\\
N 書出し D & 未達 & 未達\\
N 書出し B & 未達 & 未達\\
\hline\end{tabular}\end{center}
Gの写真FとNの写真Dだけが未達から成功へ変わり，それぞれ元GのD，元NのBと最終13保存ファイルも一致した．非supervisor保存ファイルは全ケースで介入前後不変だった．残り７状態は通知後も未達で，完了宣言を一律に足す根拠にはならない．
Gの診断は346 API・４評価・23.927秒（評価8.541秒），Nは702再送＋７通知＝709 API・14評価・64.318秒（評価19.503秒）だった．両診断とも生成０で，元の成否と不明費用は変更しない．Nの不明２生成は最後の未実行応答に限り，全API作用が既知であることを別に照合した．
同じ既知写真の事後診断であり，独立した別課題の追認ではない．Gの写真では分類に必要な現在年の成功した公開時計応答がなく，通知による公式判定の変化は公開証拠だけで要求を全て検証できたことを意味しない．
